\documentclass[11pt]{article}
\usepackage[T1]{fontenc}
\usepackage{textcomp}
\usepackage[margin=0.75in]{geometry}
\usepackage{amsmath,amssymb,amsthm}
\usepackage{array,booktabs}
\usepackage{hyperref}
\setlength{\parindent}{0pt}
\newtheorem{theorem}{Theorem}
\theoremstyle{definition}
\newtheorem{definition}{Definition}
\title{Flow Proof Artifact: Rat-mul-assoc}
\author{Generated by \texttt{flow doc proof}}
\date{Flow Proof Book}
\begin{document}
\maketitle
Rational multiplication associates.\\[0.5em]
\textbf{Source.} landau — *Foundations of Analysis*\\[0.5em]
\bigskip
\noindent{\large \textbf{Derived fact 1.} \textit{(p * q) * r = p * (q * r)} \hfill Rat \cdot multiplication \cdot parentheses do not matter}\par
\smallskip\noindent\textit{Coordinate: Rat \cdot multiplication \cdot parentheses do not matter}\par
\smallskip\noindent\textit{Source: landau}\par
\smallskip\noindent\textit{Built on: one is the left identity, for multiplication on Rat, one is the right identity, for multiplication on Rat, order does not matter, for multiplication on Rat}\par
\medskip
\noindent\textbf{Goal.} (p * q) * r = p * (q * r)\par
\noindent$\displaystyle \forall p \in Rat \forall q \in Rat \forall r \in Rat\quad (p \cdot q) * r = p * (q \cdot r)$\par
\medskip
\noindent
\begin{tabular}{@{} >{\raggedright\arraybackslash}p{0.06\textwidth} >{\raggedright\arraybackslash}p{0.47\textwidth} >{\raggedleft\arraybackslash}p{0.41\textwidth} @{}}
\textbf{\#} & \textbf{Proof} & \textbf{Mathematics} \\
\midrule
\textbf{1.} & We prove that parentheses do not matter for multiplication on Rat. & \\[0.45em]
\textbf{2.} & We invoke the definitional clause governing multiplication on Rat: one is the left identity, for multiplication on Rat (instantiated for p). & \\[0.45em]
\textbf{3.} & We invoke the derived fact governing multiplication on Rat: one is the right identity, for multiplication on Rat (instantiated for r). & \\[0.45em]
\textbf{4.} & We invoke the derived fact governing multiplication on Rat: order does not matter, for multiplication on Rat (instantiated for p, q). & \\[0.45em]
\textbf{5.} & We invoke the derived fact governing multiplication on Rat: order does not matter, for multiplication on Rat (instantiated for q, r). & \\[0.45em]
\textbf{6.} & From step 2, step 3, step 4, and step 5, this implies (p times q) times r equals p times (q times r). Hence proven. & $\displaystyle (p \cdot q) * r = p * (q \cdot r)$ \\[0.45em]
\bottomrule
\end{tabular}
\label{thm:Rat-multiplication-parentheses_do_not_matter}
\par\medskip\hrule
\end{document}
